\section{本讲小结}

导体问题先确定哪一部分等势、哪一部分保持净电荷、哪一部分可以与大地或电源交换电荷。金属内零场给内壁总量，唯一性决定空腔与外部场如何分离；屏蔽并不等同于所有感应电荷消失。

电像法的核心是边界等效。像应位于真实求解区域之外，同时满足球面或平面的电势条件与整体电荷约束。接地球、孤立中性球、孤立带电球、定势球的边界不同，中心像也不同。力由像场求，涉及自身诱导的能量要按逐渐建立的过程积分。

局部平面近似只控制邻近区域；有限导体的总感应量和附加净量分布仍需整体条件。每个近似结果都应检查几何尺度、远近极限以及净量积分，避免把局部解误用到整个导体。

